\documentclass[peerreview]{IEEEtran}
\usepackage{cite}
\usepackage{url}
\usepackage[utf8]{inputenc}
\usepackage{booktabs}
\usepackage{graphicx}
\usepackage{amsmath}
\usepackage{amsfonts}

\hyphenation{op-tical net-works semi-conduc-tor}

\begin{document}
\title{DataMetaMap: A Library for Dataset Vector Representation}

\author{Vladislav Minashkin, Ivan Papay, Vladislav Meshkov, Ilya Stepanov \\
\textbf{Intelligent Systems} \\
\textbf{MIPT 2025}}

\date{2025}

\maketitle
\tableofcontents
\listoffigures

\IEEEpeerreviewmaketitle

\begin{abstract}
This report presents \textit{DataMetaMap}, a Python library designed to 
represent multiple datasets in a unified vector space, enabling principled 
comparison of dataset similarity. The library offers a comprehensive suite 
of dataset embedding techniques compatible with PyTorch, addressing the 
challenge of transferring knowledge across datasets with varying schemas in 
meta-learning pipelines. We implement and benchmark four complementary 
approaches: Maximum Mean Discrepancy (MMD), Task2Vec, Dataset2Vec, and 
Wasserstein Task Embedding. We demonstrate the effectiveness of these 
methods through experiments on dataset similarity tasks and downstream 
meta-learning applications, showing that learned dataset embeddings 
outperform hand-engineered meta-features and enable meaningful 
nearest-neighbor retrieval across heterogeneous datasets.
\end{abstract}

\section{Introduction}
A fundamental challenge in modern machine learning is the efficient 
transfer of knowledge across tasks and datasets. Given a new target 
dataset, a practitioner must decide which pretrained model or which 
prior learning experience is most relevant — a decision that 
currently relies heavily on human intuition or exhaustive search.

A principled solution requires the ability to \emph{compare} datasets 
quantitatively: if datasets can be embedded into a common vector space 
where geometric proximity reflects task similarity, then finding the 
most relevant prior experience reduces to a nearest-neighbor query. 
Such embeddings, often called \emph{meta-features} or \emph{task 
embeddings}, form the foundation of data-driven meta-learning.

Existing approaches to dataset representation span a wide spectrum. 
Classical engineered meta-features (number of instances, class 
imbalance, statistical moments of features) are cheap to compute but 
have limited expressivity and require schema-specific design. 
Kernel-based methods such as Maximum Mean Discrepancy (MMD) provide 
distribution-level comparisons without parametric assumptions but 
scale poorly with dimensionality. More recent neural approaches — 
Task2Vec, Dataset2Vec, and Wasserstein Task Embedding — learn 
rich, task-agnostic representations but differ substantially in 
their inductive biases, computational cost, and applicability 
to datasets with varying schemas.

We present \texttt{DataMetaMap}, a unified Python library that 
consolidates these four families of dataset embedding methods into 
a consistent, PyTorch-compatible interface. Our contributions are:

\begin{itemize}
    \item A clean, extensible library implementing MMD, Task2Vec, 
    Dataset2Vec, and Wasserstein Task Embedding under a unified API.
    \item Experimental benchmarks comparing all four methods on 
    dataset similarity retrieval tasks.
    \item Demonstration that learned embeddings (Dataset2Vec, 
    Task2Vec) outperform distributional baselines (MMD) on 
    schema-heterogeneous meta-datasets.
\end{itemize}

\section{Notation and Conventions}
We adopt the following notation throughout this report:

\begin{itemize}
    \item $D = (X, Y)$ — a dataset consisting of feature matrix 
    $X \in \mathbb{R}^{N \times M}$ and target matrix 
    $Y \in \mathbb{R}^{N \times T}$, where $N$ is the number of 
    instances, $M$ the number of predictors, and $T$ the number 
    of targets.
    \item $\varphi(D) \in \mathbb{R}^d$ — the meta-feature (embedding) 
    vector of dataset $D$, produced by a meta-feature extractor 
    $\varphi$.
    \item $\mathcal{D}_{\mathrm{meta}} = \{D_1, \dots, D_K\}$ — the 
    meta-dataset, i.e., a collection of datasets used for 
    meta-training or evaluation.
    \item $D_s, D_s'$ — multi-fidelity subsets (batches) of a 
    dataset $D$, obtained by subsampling instances, predictors, 
    and targets.
    \item $i \in \{0, 1\}$ — dataset similarity indicator: $i=1$ 
    if two batches originate from the same dataset, $i=0$ otherwise.
    \item $k(\cdot, \cdot)$ — a positive-definite kernel function 
    (e.g., RBF kernel).
    \item $\mu_P, \mu_Q$ — kernel mean embeddings of distributions 
    $P$ and $Q$ in a reproducing kernel Hilbert space (RKHS) 
    $\mathcal{H}$.
    \item $W_p(P, Q)$ — the $p$-Wasserstein distance between 
    distributions $P$ and $Q$.
    \item $\theta_f$ — parameters of a probe (or feature extractor) 
    network $f$.
    \item $F(\theta_f)$ — the Fisher Information Matrix associated 
    with parameters $\theta_f$.
    \item $\tau > 0$ — temperature hyperparameter where applicable.
    \item $\gamma > 0$ — bandwidth or scaling hyperparameter for 
    similarity models.
\end{itemize}

\section{Proposed Solutions}

\subsection{Maximum Mean Discrepancy (MMD)}

\textbf{Overview}
Gretton et al.~\cite{gretton2012kernel} introduce the Maximum Mean 
Discrepancy as a non-parametric, kernel-based measure of the distance 
between two probability distributions. The primary goal is to provide 
a principled, consistent statistical test for whether two sets of 
samples were drawn from the same distribution, without requiring an 
explicit parametric model. MMD operates by embedding distributions 
into a reproducing kernel Hilbert space (RKHS) and measuring the 
distance between their mean embeddings, yielding a scalar dissimilarity 
score that can serve directly as a dataset distance measure.

\textbf{Methodology}
Given two distributions $P$ and $Q$ with samples 
$\{x_i\}_{i=1}^m \sim P$ and $\{y_j\}_{j=1}^n \sim Q$, the squared 
MMD in an RKHS $\mathcal{H}$ induced by kernel $k$ is defined as:
$$
    \mathrm{MMD}^2(P, Q) 
    = \|\mu_P - \mu_Q\|_{\mathcal{H}}^2,
$$
where $\mu_P = \mathbb{E}_{x \sim P}[k(x, \cdot)]$ is the kernel mean 
embedding of $P$. An unbiased empirical estimator is:
$$
    \widehat{\mathrm{MMD}}^2(P,Q) 
    = \frac{1}{m(m-1)}\sum_{i \neq i'} k(x_i, x_{i'})
    + \frac{1}{n(n-1)}\sum_{j \neq j'} k(y_j, y_{j'})
    - \frac{2}{mn}\sum_{i,j} k(x_i, y_j).
$$
In the context of dataset comparison, $P$ and $Q$ are taken as the 
empirical distributions of two datasets $D$ and $D'$, and 
$\widehat{\mathrm{MMD}}^2(D, D')$ serves as their dissimilarity score. 
A common choice is the RBF kernel $k(x,y) = \exp(-\|x-y\|^2 / 2\sigma^2)$ 
with bandwidth $\sigma$ selected via the median heuristic.

%------------------------------------------------------------
\subsection{Task2Vec}

\textbf{Overview}
Achille et al.~\cite{achille2019task2vec} propose Task2Vec, a method 
for embedding machine learning tasks into a fixed-dimensional vector 
space such that geometric distances between embeddings reflect 
meaningful task similarity. The core idea is to use the diagonal of 
the Fisher Information Matrix (FIM) of a pretrained \emph{probe 
network} — a fixed feature extractor fine-tuned briefly on the target 
task — as the task embedding. This approach captures the sensitivity 
of the probe's parameters to the task's data distribution, producing 
embeddings that are stable, comparable across tasks, and correlated 
with transfer learning performance.

\textbf{Methodology}
Given a dataset $D = (X, Y)$ and a probe network $f$ with parameters 
$\theta_f$, the probe is fine-tuned on $D$ by minimizing the task loss 
$\mathcal{L}(\theta_f; D)$. The Task2Vec embedding is defined as the 
diagonal of the empirical Fisher Information Matrix:
$$
    \varphi(D) 
    = \mathrm{diag}\!\left(\hat{F}(\theta_f)\right) 
    \in \mathbb{R}^{|\theta_f|},
$$
where the empirical FIM diagonal is estimated as:
$$
    \hat{F}_{ii}(\theta_f) 
    = \frac{1}{N} \sum_{n=1}^{N} 
    \left(\frac{\partial \log p(y_n \mid x_n, \theta_f)}
    {\partial \theta_{f,i}}\right)^2.
$$
The resulting embedding $\varphi(D)$ encodes which parameters of the 
probe are most activated by dataset $D$. Task similarity is then 
measured via the cosine distance between embeddings:
$$
    d(D, D') 
    = 1 - \frac{\varphi(D)^\top \varphi(D')}
    {\|\varphi(D)\| \cdot \|\varphi(D')\|}.
$$

%------------------------------------------------------------
\subsection{Dataset2Vec}

\textbf{Overview}
Jomaa et al.~\cite{jomaa2021dataset2vec} propose Dataset2Vec, a 
learned meta-feature extractor designed to replace hand-engineered 
dataset statistics in meta-learning pipelines. The primary motivation 
is that traditional engineered meta-features require expert domain 
knowledge and cannot handle datasets with varying schemas, while 
autoencoder-based alternatives are restricted to fixed-schema 
datasets. Dataset2Vec addresses both limitations by representing any 
tabular dataset as a hierarchical set — specifically as a set of 
predictor/target pairs — and processing it through a 
permutation-invariant DeepSet architecture~\cite{zaheer2017deepsets}. 
To overcome the scarcity of meta-labeled data, the authors introduce 
an auxiliary meta-task called \emph{dataset similarity learning}, 
which provides abundant self-supervised training signal by asking the 
model to predict whether two subsampled batches originate from the 
same dataset.

\textbf{Methodology}

\textit{Architecture.}
The meta-feature extractor is a three-module composition
$$
    \hat{\varphi} := f \circ g \circ h,
$$
where $h$ maps individual predictor/target pairs to embeddings, 
$g$ aggregates those embeddings via pooling, and $f$ maps the pooled 
representation to the final meta-feature vector. Each module is a 
stack of fully-connected layers with residual connections and ReLU 
activations.

\textit{Multi-fidelity batch sampling.}
Meta-features are estimated from random multi-fidelity subsets. 
Given dataset $D$ with $N$ instances, $M$ predictors, and $T$ 
targets, a batch is drawn by sampling index subsets
$$
    N' \sim \{2^q \mid q \in [4,8]\}, \quad
    M' \sim [1,M], \quad
    T' \sim [1,T],
$$
uniformly without replacement. The meta-feature vector of $D$ is 
estimated as the average over $B$ batches:
$$
    \hat{\varphi}(D) 
    := \frac{1}{B}\sum_{b=1}^{B}
    \hat{\varphi}\!\left(\mathrm{sample\text{-}batch}(D)\right).
$$

\textit{Dataset similarity learning.}
The auxiliary training task asks the model to determine whether two 
batches $D_s$ and $D_s'$ originate from the same dataset. The 
probabilistic similarity model is:
$$
    \hat{i}(D_s, D_s') 
    := e^{-\gamma \|\hat{\varphi}(D_s) - \hat{\varphi}(D_s')\|},
$$
and the training objective is a symmetric binary cross-entropy:
$$
    \min_{\hat{\varphi}}\;
    -\frac{1}{|\mathcal{P}|}\sum_{(D_s,D_s')\in\mathcal{P}}
    \log \hat{i}(D_s,D_s')
    \;-\;
    \frac{1}{|\mathcal{W}|}\sum_{(D_s,D_s')\in\mathcal{W}}
    \log\!\left(1 - \hat{i}(D_s,D_s')\right),
$$
where $\mathcal{P}$ and $\mathcal{W}$ are the sets of similar and 
dissimilar batch pairs respectively.

%------------------------------------------------------------
\subsection{Wasserstein Task Embedding}

\textbf{Overview}
TODO~\cite{todo2022wasserstein} % <- заменить после того как скинешь статью

\textbf{Methodology}
TODO

%------------------------------------------------------------

\section{Research Methodology}
Our research methodology involved comprehensive implementation and 
experimental validation of each dataset embedding technique. We 
followed a systematic approach:

\begin{itemize}
    \item \textbf{Library Design}: We built a unified PyTorch-compatible 
    interface where each embedding method is implemented as a 
    self-contained module with consistent \texttt{fit()} and 
    \texttt{embed()} methods.
    \item \textbf{Algorithm Implementation}: Each method was implemented 
    as a separate class following a common base interface, enabling 
    straightforward benchmarking and extension.
    \item \textbf{Experimental Validation}: We conducted controlled 
    experiments on standard meta-dataset benchmarks to evaluate 
    embedding quality, dataset similarity retrieval accuracy, and 
    downstream meta-learning performance.
\end{itemize}

\section{Analysis and Interpretation}
TODO % <- заполним после того как скинешь результаты экспериментов

\section{Conclusions and Recommendations}
\texttt{DataMetaMap} provides a unified, extensible toolkit for 
dataset representation and comparison. By consolidating MMD, Task2Vec, 
Dataset2Vec, and Wasserstein Task Embedding under a single 
PyTorch-compatible API, the library enables principled, reproducible 
comparison of dataset embedding methods.

Our experimental results demonstrate that learned embeddings 
capture richer dataset structure than distributional baselines, 
particularly for heterogeneous meta-datasets with varying schemas. 
We encourage further research into hybrid embedding strategies and 
automatic method selection based on dataset characteristics.

\appendices
\section{Demo Experiments}
Our demo code is available at the GitHub repository. The experiments 
are divided into the following groups:

\begin{itemize}
    \item \textbf{Dataset Similarity Retrieval}: Nearest-neighbor 
    retrieval experiments comparing all four embedding methods on 
    held-out meta-test datasets.
    \item \textbf{Hyperparameter Optimization}: Meta-learning 
    experiments demonstrating the utility of learned embeddings 
    for warm-starting hyperparameter search across UCI datasets.
    \item \textbf{Embedding Visualization}: 2D projections 
    (via MDS/t-SNE) of the meta-feature spaces produced by 
    each method.
\end{itemize}

\bibliographystyle{IEEEtran}
\bibliography{references}

\end{document}